\documentclass[a4paper,11pt]{article}
\usepackage[utf8]{inputenc}
\usepackage[english]{babel}
\usepackage{amsmath}
\usepackage{amsfonts}
\usepackage{booktabs}
\usepackage{geometry}
\usepackage{xcolor}
\usepackage{graphicx}
\usepackage{float}

\geometry{margin=2.5cm}

\title{Mathematical Analysis of Approximate Multiplication Processes}
\author{Nico Teufel}
\date{Jan 11, 2025}

\begin{document}

\maketitle

\section{Introduction}
The present system implements an approximate floating-point multiplier (AFPM) for 32-bit IEEE 754 formats. The focus is on increasing efficiency through targeted error injection in the mantissa multiplication (24-bit), controlled by a genetic structure, the so-called \textit{chromosome}.

\section{Approximation Process}

\subsection{Mantissa Decomposition (AIM\_24)}
The 24-bit mantissas are divided into three 8-bit segments: \textit{Low} (L), \textit{Medium} (M), and \textit{High} (H). The multiplication is performed by summing up 9 partial products ($mul_1$ to $mul_9$):
\begin{equation}
  Result = \sum_{i=1}^{9} mul_i \cdot 2^{k_i}
\end{equation}
Each partial product is controlled by a gene in the chromosome, which determines which 8x8-bit multiplier type (\textit{PRIM8}) is used.

\subsection{Importance Order of Partial Products}
The analysis shows a clear prioritization based on bit significance:
\begin{itemize}
  \item \textbf{Low Relevance ($2^0, 2^8$):} $mul_1$ (LL), $mul_2$ (ML), $mul_3$ (LM).
  \item \textbf{Medium Relevance ($2^{16}$):} $mul_4$ (HL), $mul_5$ (MM), $mul_6$ (LH).
  \item \textbf{High Relevance ($2^{24}, 2^{32}$):} $mul_7$ (HM), $mul_8$ (MH), $mul_9$ (HH).
\end{itemize}

\subsection{PRIM8\_xy and GroupX Approximation}
A \textit{PRIM8\_xy} multiplier decomposes the second 8-bit operand into two 4-bit nibbles.
\begin{itemize}
  \item \textbf{Index x:} Determines the approximation of the \textit{lower} nibble (bits 0-3).
  \item \textbf{Index y:} Determines the approximation of the \textit{upper} nibble (bits 4-7).
\end{itemize}
The \textit{GroupX(z)} cores inject errors by ignoring carries in the lowest \textit{z} bit positions of the column addition or replacing them with OR operations. A higher index implies a more aggressive approximation.

\section{Analysis of the 51 Chromosome Configurations}

In the file \texttt{config\_primr8.yaml}, a total of 50 active configurations were identified, which, together with the exact reference, form the 51 examined scenarios.

\subsection{Configuration Scheme}
All 50 approximate chromosomes follow a strict safety pattern to ensure numerical stability:
\begin{equation}
  Chromosom = [100, 100, 100, G, G, G, 0, 0, 0]
\end{equation}
\begin{itemize}
  \item \textbf{Genes 0-2 (Low):} Value \texttt{100} (\texttt{zero\_mul}). These partial products are completely ignored.
  \item \textbf{Genes 3-5 (Medium):} Value \texttt{G}. Here the specific approximation type $P_{xy}$ is applied.
  \item \textbf{Genes 6-8 (High):} Value \texttt{0} (\texttt{PRIM8\_11}). These critical products are calculated exactly.
\end{itemize}

\subsection{The Pxy Families}
The 49 examined $P_{xy}$ configurations (plus the exact one) are divided as follows:
\begin{table}[h]
  \centering
  \begin{tabular}{lll}
    \toprule
    \textbf{Family} & \textbf{Range y} & \textbf{Description}                 \\
    \midrule
    P41 -- P44      & $y \in [1, 4]$   & Moderate Approximation (x=4)         \\
    P51 -- P55      & $y \in [1, 5]$   & ...                                  \\
    P61 -- P66      & $y \in [1, 6]$   & ...                                  \\
    P71 -- P77      & $y \in [1, 7]$   & ...                                  \\
    P81 -- P88      & $y \in [1, 8]$   & ...                                  \\
    P91 -- P99      & $y \in [1, 9]$   & Strong Approximation                 \\
    Pa1 -- Paa      & $y \in [1, 10]$  & Maximum Approximation (Hex 'a' = 10) \\
    \bottomrule
  \end{tabular}
  \caption{Overview of the 49 approximate configurations.}
\end{table}

\section{Visual Analysis of Errors}
For better illustration of the error behavior, the histograms for selected configurations are analyzed.

The histograms were generated using a statistical simulation with $N = 2 \cdot 10^6$ samples. For each configuration, pairs of floating-point numbers $(A, B)$ were drawn from a uniform distribution in the range $[-1000, 1000]$. The relative error $e_{rel}$ was calculated as:
\begin{equation}
  e_{rel} = \frac{A_{approx} - A_{exact}}{A_{exact}}
\end{equation}
where $A_{exact} = A \cdot B$ is the exact product and $A_{approx} = \text{AFPM}(A, B)$ is the result of the approximate multiplication. The plots show the density distribution of $e_{rel}$ in percent. Additionally, the Mean Bias ($\mu = \frac{1}{N}\sum e_{rel}$) and the Standard Deviation ($\sigma$) are indicated.

The results are grouped into families according to the index $x$ (e.g., Family P4, P5, etc.). This grouping is appropriate as the $x$ parameter in the $P_{xy}$ approximation represents the most important factor for the error characteristics, while the $y$ parameter allows for secondary fine-tuning within each family.

Below are the histograms of the relative error distributions for all examined configuration families (P4x to Pax).

\subsection{Placeholder for Visualizations}
% \subsubsection{Family P4}
% \begin{figure}[H]
%   \centering
%   \begin{minipage}{0.48\textwidth}
%     \centering
%     \includegraphics[width=\linewidth]{experiments/hist_P41.png}
%     \caption*{P41}
%   \end{minipage}
%   \hfill
%   \begin{minipage}{0.48\textwidth}
%     \centering
%     \includegraphics[width=\linewidth]{experiments/hist_P42.png}
%     \caption*{P42}
%   \end{minipage}
%   \hfill
%   \begin{minipage}{0.48\textwidth}
%     \centering
%     \includegraphics[width=\linewidth]{experiments/hist_P43.png}
%     \caption*{P43}
%   \end{minipage}
%   \hfill
%   \begin{minipage}{0.48\textwidth}
%     \centering
%     \includegraphics[width=\linewidth]{experiments/hist_P44.png}
%     \caption*{P44}
%   \end{minipage}
%   \hfill
%   \caption{Histograms of family P4}
% \end{figure}
% \clearpage

% \subsubsection{Family P5}
% \begin{figure}[H]
%   \centering
%   \begin{minipage}{0.48\textwidth}
%     \centering
%     \includegraphics[width=\linewidth]{experiments/hist_P51.png}
%     \caption*{P51}
%   \end{minipage}
%   \hfill
%   \begin{minipage}{0.48\textwidth}
%     \centering
%     \includegraphics[width=\linewidth]{experiments/hist_P52.png}
%     \caption*{P52}
%   \end{minipage}
%   \hfill
%   \begin{minipage}{0.48\textwidth}
%     \centering
%     \includegraphics[width=\linewidth]{experiments/hist_P53.png}
%     \caption*{P53}
%   \end{minipage}
%   \hfill
%   \begin{minipage}{0.48\textwidth}
%     \centering
%     \includegraphics[width=\linewidth]{experiments/hist_P54.png}
%     \caption*{P54}
%   \end{minipage}
%   \hfill
%   \begin{minipage}{0.48\textwidth}
%     \centering
%     \includegraphics[width=\linewidth]{experiments/hist_P55.png}
%     \caption*{P55}
%   \end{minipage}
%   \hfill
%   \caption{Histograms of family P5}
% \end{figure}
% \clearpage

% \subsubsection{Family P6}
% \begin{figure}[H]
%   \centering
%   \begin{minipage}{0.48\textwidth}
%     \centering
%     \includegraphics[width=\linewidth]{experiments/hist_P61.png}
%     \caption*{P61}
%   \end{minipage}
%   \hfill
%   \begin{minipage}{0.48\textwidth}
%     \centering
%     \includegraphics[width=\linewidth]{experiments/hist_P62.png}
%     \caption*{P62}
%   \end{minipage}
%   \hfill
%   \begin{minipage}{0.48\textwidth}
%     \centering
%     \includegraphics[width=\linewidth]{experiments/hist_P63.png}
%     \caption*{P63}
%   \end{minipage}
%   \hfill
%   \begin{minipage}{0.48\textwidth}
%     \centering
%     \includegraphics[width=\linewidth]{experiments/hist_P64.png}
%     \caption*{P64}
%   \end{minipage}
%   \hfill
%   \begin{minipage}{0.48\textwidth}
%     \centering
%     \includegraphics[width=\linewidth]{experiments/hist_P65.png}
%     \caption*{P65}
%   \end{minipage}
%   \hfill
%   \begin{minipage}{0.48\textwidth}
%     \centering
%     \includegraphics[width=\linewidth]{experiments/hist_P66.png}
%     \caption*{P66}
%   \end{minipage}
%   \hfill
%   \caption{Histograms of family P6}
% \end{figure}
% \clearpage

% \subsubsection{Family P7}
% \begin{figure}[H]
%   \centering
%   \begin{minipage}{0.48\textwidth}
%     \centering
%     \includegraphics[width=\linewidth]{experiments/hist_P71.png}
%     \caption*{P71}
%   \end{minipage}
%   \hfill
%   \begin{minipage}{0.48\textwidth}
%     \centering
%     \includegraphics[width=\linewidth]{experiments/hist_P72.png}
%     \caption*{P72}
%   \end{minipage}
%   \hfill
%   \begin{minipage}{0.48\textwidth}
%     \centering
%     \includegraphics[width=\linewidth]{experiments/hist_P73.png}
%     \caption*{P73}
%   \end{minipage}
%   \hfill
%   \begin{minipage}{0.48\textwidth}
%     \centering
%     \includegraphics[width=\linewidth]{experiments/hist_P74.png}
%     \caption*{P74}
%   \end{minipage}
%   \hfill
%   \begin{minipage}{0.48\textwidth}
%     \centering
%     \includegraphics[width=\linewidth]{experiments/hist_P75.png}
%     \caption*{P75}
%   \end{minipage}
%   \hfill
%   \begin{minipage}{0.48\textwidth}
%     \centering
%     \includegraphics[width=\linewidth]{experiments/hist_P76.png}
%     \caption*{P76}
%   \end{minipage}
%   \hfill
% \end{figure}
% \begin{figure}[H]
%   \centering
%   \begin{minipage}{0.48\textwidth}
%     \centering
%     \includegraphics[width=\linewidth]{experiments/hist_P77.png}
%     \caption*{P77}
%   \end{minipage}
%   \hfill
%   \caption{Histograms of family P7}
% \end{figure}
% \clearpage

% \subsubsection{Family P8}
% \begin{figure}[H]
%   \centering
%   \begin{minipage}{0.48\textwidth}
%     \centering
%     \includegraphics[width=\linewidth]{experiments/hist_P81.png}
%     \caption*{P81}
%   \end{minipage}
%   \hfill
%   \begin{minipage}{0.48\textwidth}
%     \centering
%     \includegraphics[width=\linewidth]{experiments/hist_P82.png}
%     \caption*{P82}
%   \end{minipage}
%   \hfill
%   \begin{minipage}{0.48\textwidth}
%     \centering
%     \includegraphics[width=\linewidth]{experiments/hist_P83.png}
%     \caption*{P83}
%   \end{minipage}
%   \hfill
%   \begin{minipage}{0.48\textwidth}
%     \centering
%     \includegraphics[width=\linewidth]{experiments/hist_P84.png}
%     \caption*{P84}
%   \end{minipage}
%   \hfill
%   \begin{minipage}{0.48\textwidth}
%     \centering
%     \includegraphics[width=\linewidth]{experiments/hist_P85.png}
%     \caption*{P85}
%   \end{minipage}
%   \hfill
%   \begin{minipage}{0.48\textwidth}
%     \centering
%     \includegraphics[width=\linewidth]{experiments/hist_P86.png}
%     \caption*{P86}
%   \end{minipage}
%   \hfill
% \end{figure}
% \begin{figure}[H]
%   \centering
%   \begin{minipage}{0.48\textwidth}
%     \centering
%     \includegraphics[width=\linewidth]{experiments/hist_P87.png}
%     \caption*{P87}
%   \end{minipage}
%   \hfill
%   \begin{minipage}{0.48\textwidth}
%     \centering
%     \includegraphics[width=\linewidth]{experiments/hist_P88.png}
%     \caption*{P88}
%   \end{minipage}
%   \hfill
%   \caption{Histograms of family P8}
% \end{figure}
% \clearpage

% \subsubsection{Family P9}
% \begin{figure}[H]
%   \centering
%   \begin{minipage}{0.48\textwidth}
%     \centering
%     \includegraphics[width=\linewidth]{experiments/hist_P91.png}
%     \caption*{P91}
%   \end{minipage}
%   \hfill
%   \begin{minipage}{0.48\textwidth}
%     \centering
%     \includegraphics[width=\linewidth]{experiments/hist_P92.png}
%     \caption*{P92}
%   \end{minipage}
%   \hfill
%   \begin{minipage}{0.48\textwidth}
%     \centering
%     \includegraphics[width=\linewidth]{experiments/hist_P93.png}
%     \caption*{P93}
%   \end{minipage}
%   \hfill
%   \begin{minipage}{0.48\textwidth}
%     \centering
%     \includegraphics[width=\linewidth]{experiments/hist_P94.png}
%     \caption*{P94}
%   \end{minipage}
%   \hfill
%   \begin{minipage}{0.48\textwidth}
%     \centering
%     \includegraphics[width=\linewidth]{experiments/hist_P95.png}
%     \caption*{P95}
%   \end{minipage}
%   \hfill
%   \begin{minipage}{0.48\textwidth}
%     \centering
%     \includegraphics[width=\linewidth]{experiments/hist_P96.png}
%     \caption*{P96}
%   \end{minipage}
%   \hfill
% \end{figure}
% \begin{figure}[H]
%   \centering
%   \begin{minipage}{0.48\textwidth}
%     \centering
%     \includegraphics[width=\linewidth]{experiments/hist_P97.png}
%     \caption*{P97}
%   \end{minipage}
%   \hfill
%   \begin{minipage}{0.48\textwidth}
%     \centering
%     \includegraphics[width=\linewidth]{experiments/hist_P98.png}
%     \caption*{P98}
%   \end{minipage}
%   \hfill
%   \begin{minipage}{0.48\textwidth}
%     \centering
%     \includegraphics[width=\linewidth]{experiments/hist_P99.png}
%     \caption*{P99}
%   \end{minipage}
%   \hfill
%   \caption{Histograms of family P9}
% \end{figure}
% \clearpage

% \subsubsection{Family Pa}
% \begin{figure}[H]
%   \centering
%   \begin{minipage}{0.48\textwidth}
%     \centering
%     \includegraphics[width=\linewidth]{experiments/hist_Pa1.png}
%     \caption*{Pa1}
%   \end{minipage}
%   \hfill
%   \begin{minipage}{0.48\textwidth}
%     \centering
%     \includegraphics[width=\linewidth]{experiments/hist_Pa2.png}
%     \caption*{Pa2}
%   \end{minipage}
%   \hfill
%   \begin{minipage}{0.48\textwidth}
%     \centering
%     \includegraphics[width=\linewidth]{experiments/hist_Pa3.png}
%     \caption*{Pa3}
%   \end{minipage}
%   \hfill
%   \begin{minipage}{0.48\textwidth}
%     \centering
%     \includegraphics[width=\linewidth]{experiments/hist_Pa4.png}
%     \caption*{Pa4}
%   \end{minipage}
%   \hfill
%   \begin{minipage}{0.48\textwidth}
%     \centering
%     \includegraphics[width=\linewidth]{experiments/hist_Pa5.png}
%     \caption*{Pa5}
%   \end{minipage}
%   \hfill
%   \begin{minipage}{0.48\textwidth}
%     \centering
%     \includegraphics[width=\linewidth]{experiments/hist_Pa6.png}
%     \caption*{Pa6}
%   \end{minipage}
%   \hfill
% \end{figure}
% \begin{figure}[H]
%   \centering
%   \begin{minipage}{0.48\textwidth}
%     \centering
%     \includegraphics[width=\linewidth]{experiments/hist_Pa7.png}
%     \caption*{Pa7}
%   \end{minipage}
%   \hfill
%   \begin{minipage}{0.48\textwidth}
%     \centering
%     \includegraphics[width=\linewidth]{experiments/hist_Pa8.png}
%     \caption*{Pa8}
%   \end{minipage}
%   \hfill
%   \begin{minipage}{0.48\textwidth}
%     \centering
%     \includegraphics[width=\linewidth]{experiments/hist_Pa9.png}
%     \caption*{Pa9}
%   \end{minipage}
%   \hfill
%   \begin{minipage}{0.48\textwidth}
%     \centering
%     \includegraphics[width=\linewidth]{experiments/hist_Paa.png}
%     \caption*{Paa}
%   \end{minipage}
%   \hfill
%   \caption{Histograms of family Pa}
% \end{figure}
% \clearpage

\section{Further Research}
This analysis enables further research. A few areas are presented below:
\begin{itemize}
  \item What do the histograms for matrix multiplications look like?
  \item How can a bias correction be applied without drawbacks on the efficiency of the approximation process?
  \item Does the iterative refinement process for the 1D Poisson Problem converge (with or without bias correction)?
  \item Can we model the approximate multipliers mathematically and derive properties and bounds from this model?
\end{itemize}

\end{document}